\section{Adaptive, redundancy-aware probe generation}
\label{sec:probes}

\subsection{The redundancy problem}
The full coalition family (mean and maximum over every non-trivial subset of
criteria) has size exponential in $m$. Worse, real criteria are often
\emph{redundant}: several measured objectives track the same underlying concern
and are strongly positively correlated. Averaging selection rules implicitly
\emph{over-weight} such a group---it is counted once per member---biasing the
recommendation. 

Consider a simple illustrative example: a 3-criteria problem where $f_1$ is cost, and $f_2$ and $f_3$ are two identical measures of environmental impact. An averaging rule will implicitly double-weight the environmental impact simply because it was measured twice. We need a probe family that (i) is small, (ii) stays monotone,
and (iii) is less sensitive to redundancy than averaging aggregation. LUR with adaptive probes identifies that $f_2$ and $f_3$ are identical (correlation 1.0) and clusters them, generating a reduced set of probes that does not repeatedly average the redundant pair.

\subsection{Correlation clustering}
We group \emph{redundant} criteria by single-linkage connected components of the
graph whose edges join pairs with Pearson correlation $\Corr(f_i,f_j)\ge\theta$. We chose single-linkage clustering intentionally because it is conservative regarding redundancy. Two criteria only need a single path of high correlations between them to be grouped together. This prevents the artificial slicing of a continuous underlying dimension into multiple independent criteria that would otherwise overwhelm the probe generation.

\emph{Redundancy is positive correlation:} two criteria that rise and fall
together carry duplicate information, whereas negatively correlated criteria are
genuine trade-offs and must remain separate---hence we threshold the signed
correlation, not its absolute value. This yields clusters $C_1,\dots,C_c$.

\subsection{The adaptive probe family}
LUR always includes the canonical anchors---the $m$ singletons, the grand mean,
and the grand maximum---and adds one mean and one maximum probe \emph{per cluster}
of size $>1$:
\begin{equation}
\Q=\{q_i\}_{i=1}^m\;\cup\;\{q_\mu,q_{\max}\}\;\cup\;
\bigcup_{|C_k|>1}\{q_{C_k,\mathrm{mean}},\,q_{C_k,\max}\}.
\end{equation}
A redundant group thus contributes \emph{one} aggregate (cluster) question rather
than a separate aggregate per member, so the family does not amplify a
correlated group the way an averaging aggregate does, while the count stays
$O(m+c)$---linear, not exponential. We stress that this \emph{reduces}, but does
not eliminate, redundancy sensitivity: the singleton probes for the group's
members remain in $\Q$ (they are needed for separation, Theorem~\ref{thm:pareto}),
so a redundant group can still influence lower-order lexicographic tie-breaking.
The claim is therefore empirical---less redundancy bias than averaging methods on
the tested benchmarks (\S\ref{sec:exp})---not a duplication-invariance theorem.
All probes are non-negative-weighted sums or maxima, hence monotone, so
Pareto compatibility (Theorem~\ref{thm:pareto}) is preserved.

\begin{theorem}[Non-negative-probe Pareto compatibility]
\label{thm:nonneg}
If every probe in $\Q$ is a non-negative-weighted sum or a maximum of the
normalised criteria, then LUR is Pareto compatible.
\end{theorem}

\subsection{A perturbation bound (not a ranking guarantee)}
We state the reduction effect as a regret-value perturbation lemma with explicit
constants. We are deliberately modest: it bounds the change in individual
disappointments, \emph{not} the change in the LUR ranking or winner, and it
depends on the normalisation floor $\varepsilon$, which is decision-critical.

\begin{lemma}[Disappointment perturbation]
\label{thm:approx}
Suppose for a full-family probe $q$ there is a reduced-family probe $q'$ with
$\sup_r|q(r)-q'(r)|\le\eta$ (hence also $|q^\star-q'^\star|\le\eta$ and
$|q^{\mathrm w}-q'^{\mathrm w}|\le\eta$), and let the normalisation ranges satisfy
$q^{\mathrm w}-q^\star,\,q'^{\mathrm w}-q'^\star\ge\beta>0$. Then for every $x$
\[
|\reg_q(x)-\reg_{q'}(x)|\;\le\;\frac{3\eta}{\beta},
\]
with the floor $\beta\ge\varepsilon$ in the worst case. The bound is vacuous when
$\beta$ is of order $\eta$, and informative only when the probe's attainable
range is bounded away from zero---a regime the analyst can check directly.
\end{lemma}

Whether the reduced family admits such an $\eta$-close $q'$ for every full-family
probe is \emph{not} established by correlation clustering in general; we therefore
present Lemma~\ref{thm:approx} as a sensitivity statement and rely on the
empirical redundancy study (\S\ref{sec:exp}) for evidence that clustering retains
decision quality.

\subsection{Empirical reduction}
Table~\ref{tab:reduction} reports the realised reduction on concave instances.
The full family here is $\{q_{\mathrm{mean},S},q_{\max,S}:\emptyset\neq
S\subseteq\{1,\dots,m\}\}$, of size $2(2^m-1)-m$ (mean and max coincide on
singletons); the adaptive family is linear in $m$, giving a nearly four-thousand
-fold reduction by $m=15$ with, as \S\ref{sec:exp} shows, little loss of decision
quality.

\begin{table}[t]
\centering\small
\caption{Probe-family size: full coalition family ($2(2^m-1)-m$ probes) vs.\
adaptive correlation-clustering family (concave instances, $N=400$).}
\label{tab:reduction}
\begin{tabular}{rrrr}
\toprule
$m$ & full coalitions & adaptive probes & reduction\\
\midrule
3  & 11     & 5  & 2.2$\times$\\
5  & 57     & 7  & 8.1$\times$\\
8  & 502    & 10 & 50.2$\times$\\
10 & 2{,}036 & 12 & 169.7$\times$\\
15 & 65{,}519 & 17 & 3{,}854.1$\times$\\
\bottomrule
\end{tabular}
\end{table}

\subsection{Greedy probe selection (optional)}
When even the adaptive family is larger than needed, a greedy set-cover heuristic
selects a minimal subfamily distinguishing all dominant pairs; for the
max-coverage variant it attains a $(1-1/e)$ guarantee and for full coverage a
$1+\ln d$ guarantee \citep{nemhauser1978,chvatal1979}. We do not use it in the
experiments, which already run on the compact adaptive family.
